\documentclass[10pt]{article}

\usepackage[preprint]{tmlr}
\usepackage[utf8]{inputenc}
\usepackage[T1]{fontenc}
\usepackage{hyperref}
\usepackage{url}
\usepackage{booktabs}
\usepackage{amsfonts}
\usepackage{amsmath}
\usepackage{microtype}
\usepackage{graphicx}
\usepackage{xcolor}
\usepackage{underscore}
\usepackage{longtable}

\title{Evidence, Not Enrollment: Eight Recomputable Methodological\\ Diagnostics on 2026 AI Frontier Claims}

\author{%
  Anonymous Author(s) \\
  Affiliation \\
  Address \\
  \texttt{email}
}

\def\month{MM}
\def\year{YYYY}
\def\openreview{\url{https://openreview.net/forum?id=XXXX}}

\begin{document}

\maketitle

\begin{abstract}
The 2026 discussion of AI research fronts---world models, agents, embodied and
physical AI, AI for science, trustworthy AI, and foundation-model engineering---
often presents six directions as one synchronous wave of progress. We audit the
claims behind each direction and find that none of the six is wrong; the error is
placing results of \emph{different evidence grades} in the same table. We make the
audit operational with eight \emph{recomputable} methodological diagnostics---each
built to a pre-registered criterion, carried out as a three-arm comparison, a
statistical-power calculation, an arithmetic self-consistency check, a null-model
geometry computation, or a stratified re-computation of published aggregates, and
each equipped with counter-evidence self-tests that prove the apparatus can fail
and still be caught. The diagnostics deliberately target ``substrate
unavailable'' cards: when the code, weights, or per-trial logs of a method are not
released, we convert ``cannot test'' into a quantitative statement about what
conditions would make it testable, whether the published numbers are themselves
inconsistent, whether a reported paradox is the arithmetic of a low-rank subspace,
or whether an aggregate mean is driven by a few scenarios. Four diagnostics return
a directly citable verdict: an action-conditioned next-state predictor has the
lowest \emph{absolute} long-horizon error (0.0024) even though the
``lowest-slope'' criterion fails against a quantization floor; multi-view
projection degrades 3.03x under cross-embodiment perturbation versus 8.50x for a
single view; early-stopping success is governed by the detector's false-positive
rate, collapsing once FPR reaches 0.01; and an aggregate success rate of 92.1\%
is a flat result across eight scenarios, not a mean inflated by a few strong ones
(leave-one-out spread 2.38pp, chi-square p = 0.40). Across all eight, the scope of
a conclusion is fixed by the apparatus and reporting caliber, not by model scale.
Our reproducible lesson: every criterion must carry its own counter-evidence, and
unavailability is not the same as un-informative.
\end{abstract}

\section{Introduction: Naming the Shift, Auditing the Evidence}
\label{sec:intro}

``From predicting words to predicting world states'' is a popular summary of where
2026 AI research is concentrating. The academic term is \emph{next-state
prediction} (NSP). Two things follow once we check the underlying claims rather
than the framing. First, NSP is not a single invention but a continuation of a
line running from at least 2018~\cite{worldmodels2018} to late-2026 world-action
models~\cite{wam2026}. Second, treating NSP as a unified objective that
\emph{replaces} next-token prediction has, at the time of writing, no
fixed-backbone three-arm comparison behind it, and there is direct evidence that
state tracking is a distributed by-product of next-token
heads~\cite{statelaws2025}. This paper is organized around that gap: separating
the narrative from the evidence.

The broader motivation is that the six directions---world models, agents, embodied
and physical AI, AI for science, trustworthy AI, and foundation-model
engineering---are frequently narrated as one synchronized surge. A claim audit
quickly shows that the individual results are real but of \emph{unequal evidence
grade}: some are recomputable, many are author-reported with no released
substrate, and a few are internally inconsistent with their own public numbers.
Our contribution is to make the audit \emph{operational} rather than rhetorical.
We design eight diagnostics, each with a pre-registered criterion, a concrete
apparatus (a three-arm controlled comparison, an a-priori power calculation, an
arithmetic self-check, a null-model geometry computation, or a stratified
re-computation of published aggregates), and a set of counter-evidence self-tests
that prove the apparatus can produce a red verdict and that we can catch it.

We state the headline up front. Across the eight diagnostics, \emph{the scope of
each conclusion is fixed by the apparatus and its reporting caliber, not by model
scale}. This is the single most portable lesson in the paper, and the rest is
its demonstration. Four diagnostics return a directly citable verdict (Section
\ref{sec:diagnostics}); four return a quantitative statement about what can and
cannot be concluded from public material alone, which we treat as a result rather
than a gap. Section~\ref{sec:scope} states the common methodological skeleton and
the evidence taxonomy that motivates the whole exercise; Section~\ref{sec:lessons}
distills the reusable lessons; Section~\ref{sec:limits} states limits and citation
discipline.
\section{Scope and Method: A Shared Skeleton for Recomputable Diagnostics}
\label{sec:scope}

\subsection{Evidence taxonomy as motivation}

Before building diagnostics we classify the claims. Each of 33 assertions drawn
from the six directions is scored on three axes: the \emph{carrier} of the result
(A1 = independent reproduction, A2 = recomputable from released data, A3 =
author-reported with partial artifact, A4 = announcement only); the \emph{number
caliber} (B1 = ratio with explicit denominator, B2 = self-reported ratio, B3 =
qualitative or undefined); and the \emph{recomputability} (C1 = fully
recomputable, C2 = recomputable with caveats, C3 = not recomputable from public
material). The composite is the geometric mean of the three axis weights. Of the
33 assertions, 20 resolve to \texttt{PERF\_EVIDENCE} and 13 to
\texttt{ANNOUNCEMENT\_ONLY}; the per-direction composite means range from 0.43
(AI for science) to 0.66 (trustworthy AI). \texttt{ANNOUNCEMENT\_ONLY} is not a
claim that a result is false---it records that the current public material does
not let anyone independently recompute it. That distinction is the reason we
build diagnostics instead of simply downgrading every number.

\subsection{The shared diagnostic skeleton}

All eight diagnostics follow the same five-step skeleton, and each step has an
explicit failure mode we guard against:

\begin{enumerate}
  \item \textbf{Pre-registered criterion.} A single, falsifiable criterion is fixed
  \emph{before} running, so the result is judged by a rule rather than by the
  outcome we would prefer. (Failure mode: a criterion whose direction is
  backwards, which we hit on the 1.9x consistency check and had to repair.)
  \item \textbf{Controlled apparatus.} A three-arm comparison with matched
  backbone/data/compute, an a-priori power calculation, an arithmetic self-check,
  a null-model geometry computation, or a stratified re-computation---whatever
  matches the substrate available.
  \item \textbf{Counter-evidence self-test.} An injected violation must turn the
  verdict red. A criterion that cannot fail is decoration, and the self-tests are
  what separate a real gate from a tautology.
  \item \textbf{Landed, recomputable evidence.} Every number is written to a file
  from which it can be re-derived, with the command line recorded.
  \item \textbf{Honest caliber statement.} The report states exactly what the
  apparatus does and does not measure, so the scope of the conclusion travels with
  the number.
\end{enumerate}

The recurring lesson, made concrete across eight instances, is that step~3 is the
one most often skipped and most often load-bearing: a quantization floor, a
mis-aligned arm, a template shortcut, a cancellation artifact, a ``nan'' in a
non-central-$t$ tail, a reversed criterion direction, and a slow run that
disagrees with a table all \emph{looked} green until a counter-evidence test
caught them. Section~\ref{sec:lessons} treats this as the central methodological
contribution.
\section{The Eight Diagnostics}
\label{sec:diagnostics}

Table~\ref{tab:eight} lists all eight, each with its apparatus, pre-registered
criterion, main result, verdict, and the caliber boundary that travels with it.
Four are direct, citable findings; four are quantitative statements about
public-material limits that we nonetheless count as results.

\begin{center}
\scriptsize
\setlength{\tabcolsep}{3.2pt}
\begin{longtable}{@{}p{2.3cm}p{2.5cm}p{2.5cm}p{3.6cm}p{1.85cm}p{2.6cm}@{}}
\toprule
Diagnostic (project / card) & Apparatus & Pre-registered criterion & Main result & Verdict & Caliber boundary \\
\midrule
\endhead
NSP three-arm (worldmodel E1) & synthetic linear dynamics, shared backbone, heads differ & state arm has lowest cumulative slope & state slope 0.197 not lowest; absolute mse@1000 = 0.0024 lowest & criterion \textbf{failed}; keep secondary finding & quantization floor fakes low slope; not a paper reproduction \\
Multi-view (worldmodel E2) & 4 complementary projections, matched params 68748 & multi cross-embodiment degradation $<$ single & 3.03x vs 8.50x, 3/3 seeds & \textbf{h2\_supported = true} & single-step, not E1 long-horizon; CF-WAM code unreleased \\
DoM probe (trustworthy E2) & Qwen2.5-7B/14B real weights, unseen-family transfer & Q1 unseen TPR $\geq$ 0.8 + Q2 signal$-$shortcut $\geq$ 0.2 & 7B Q1/Q2 marginal; 14B validator TPR = 0, Q2 = 0.10 fail & \textbf{mixed / partly negative} & cross-scale hidden 3584 vs 5120 not portable; not a cross-family reproduction \\
Termination-lag early stop (agentic E1) & pure-numpy synthetic patch trajectories, three arms & zero cancellation and success not degraded; dominant factor is FPR & calibrated cancel 3.05\%, FPR = 0 works, FPR $\geq$ 0.01 success collapses & \textbf{criterion holds} (fpr\_dominates = true) & synthetic params, not measured; FPR/FNR are a sweep grid \\
Workflow detectability (agentic E2) & a-priori non-central-$t$ power + 5000 MC & SESOI unavailable $\rightarrow$ report MDE; difficulty interaction harder than main & main MDE dz 0.234, interaction d 0.468 (2x), R16 $\rightarrow$ 0.059 & \textbf{cannot-judge is the result} (not\_answerable) & MDE is not existence; dz unitless, $\sigma$ unpublished so no score points \\
1.9x consistency (ai4science E2) & arithmetic self-check on public numbers + identifiability & headline should be reproducible from public sub-items & mean 2.142 / median 2.05 / geom 2.112 all $\neq$ 1.9; feasible interval [1.999, 2.400] excludes 1.9 & \textbf{not\_identifiable\_with\_public\_data} & does not claim 1.9 wrong, only unreproducible; direction not identifiable without pairs \\
Explanation--energy geometry (fm E2) & isotropic null model d=128/r=4 + Beta analytic & is the 89.82\% vs 9.52\% split a low-rank geometry universal? & E $\approx$ r/d = 3.125\% (median 2.63\%), corr(A,E) = 0.011, A\_null 99th = 0.220 matches AROMA 0.219 & \textbf{GEOMETRY\_GENERIC} (C1--C4 all true) & not AROMA reproduction; h isotropic is an assumption, holds also for anisotropic; +8.00pp count gain not re-measured \\
92.1\% stratified re-compute (embodied E2) & ACE-CaP-0 Table 1 per-scenario aggregates, 8 scenarios n=30 & is 92.1\% a few-scenario high mean; is the tail stable? & balanced design (weighted == unweighted 0.9208), LOO range 2.38pp, lowest scenario 83.3\% Wilson CI [66.4\%, 92.7\%] covers overall, $\chi^2(7) = 7.259$, p = 0.402, variance ratio 1.04 & \textbf{STRATIFIED\_RECOMPUTE\_PASS} (C1--C4 true) & not physical reproduction, only Table 1 aggregates, no per-trial series (between-scenario variance computable, within-scenario not); n=30 ``not significant'' $\neq$ ``proven stable''; GR00T2 is a different experiment \\
\bottomrule
\caption{The eight recomputable methodological diagnostics, with apparatus,
pre-registered criterion, main result, verdict, and the caliber boundary that
travels with each conclusion. All numbers are re-derivable from the corresponding
project's \texttt{evidence/} artifacts and \texttt{analyses/} scripts.}
\label{tab:eight}
\end{longtable}
\end{center}

\subsection{NSP three-arm (world model E1)}

On a synthetic linear dynamical system, three arms share backbone, data, and
compute and differ only in the output head (token cross-entropy over a
quantized frame, frame-space MSE, or latent-state MSE). The pre-registered
criterion was that the state arm should have the \emph{lowest cumulative
rollout slope}. It failed: the state slope 0.197 was not the lowest. The token
arm's slope (0.136) was lowest only because its starting error hit a quantization
floor, which fabricates a low slope on a stable system. The state arm does,
however, have the lowest \emph{absolute} long-horizon error (mse@1000 =
0.0024, 0.249x of frame and 0.066x of token), and an action ablation ratio of
101.9 confirms the model uses the action signal. The verdict is therefore
``criterion failed, keep a secondary finding'': the abstract level of the
prediction target affects absolute precision (support) but not cumulative slope
(no support), and the slope criterion is unreliable on stable dynamics. The
direct counterpoint to the narrative is that the paper-level claim of a
fixed-backbone three-arm advantage has no precedent in the 50 abstracts we
retrieved, and the cited negative result is that state tracking is limited to
pattern matching and emerges from weak interaction of several next-token
heads~\cite{statelaws2025}.

\subsection{Multi-view projection (world model E2)}

Because the CF-WAM code and weights are not released~\cite{wam2026}, we build a
synthetic analogue: a latent state observed through four complementary
projections, with a single-view arm (visual only) and a multi-view arm (four
views concatenated), matched at 68{,}748 parameters. Across three seeds the
multi-view arm degrades 3.03x under a cross-embodiment perturbation versus
8.50x for the single view (a 5.47 gap), with the direction stable 3/3 seeds;
criterion H2 is supported. The honest boundary is that information complementarity
is not free (the multi-view input is higher-dimensional), that robustness---not
in-distribution accuracy (0.0081 vs 0.0091, ratio 0.891)---is the real analogue
of ``cross-embodiment generalization'', and that an under-trained smoke run had
Q2 false while the full run was stable. We report that smoke instability rather
than hiding it.

\subsection{DoM probe on real models (trustworthy AI E2)}

Using Qwen2.5-7B-base and 14B-base weights, we train a probe only on a
test-suite-tampering family plus five honest templates, and test transfer to
fully unseen validator/scorefile/acquire families. The verdict is
mixed/partly negative. On 7B all three unseen families reach TPR@1\%FPR = 1.0
(Q1 passes) but a label-permutation control reaches AUC 0.758, so
signal$-$shortcut = 0.242 only barely clears the 0.2 bar (Q2 marginal). On 14B
the validator family has TPR = 0.0 (Q1 fails) and the permutation control reaches
AUC 0.85 (Q2 fails). The hidden dimensions differ (3584 vs 5120), so the
direction does not transfer across scale, and Q3 is judged
NOT\_ANSWERABLE. Logit-lens readouts are garbled tokens, not the
``cheating/illicit'' semantics of the original DoM report~\cite{dom2026}---the
probe direction is not interpretable in our setup. This is a method-level result
about probe transfer, not a reproduction of the paper.

\subsection{Termination-lag early stop (agentic E1)}

The per-step patch trajectories of the termination-lag study~\cite{termlag2026}
are not released, so we build an equivalent synthetic trajectory apparatus
(three arms: baseline, offline oracle, deployable detector with a
FPR/FNR grid). The baseline cancellation rate calibrates to 3.05\% (inside the
paper's 1--8\% band) and the offline oracle saves 70.4\% of steps with zero
cancellation; but the online-deployable detector is only effective at FPR = 0,
with success collapsing from 96.95\% to 94.35\% at FPR = 0.01 and to 62.94\% at
FPR = 0.30. The criterion ``early-stop success is governed by the
false-positive rate, not recall'' holds (fpr\_dominates = true): a false negative
only stops later, but a false positive kills an unfinished run outright. This
translates the paper's offline-oracle conclusion into an online deployment
constraint: you need a near-zero-FPR correctness criterion before early stopping,
or the saved steps cost success rate. The self-tests (FPR = 1 must judge false;
FNR = 1 must judge false; perfect detector must judge true) prove the apparatus
can fail.

\subsection{Workflow-author detectability (agentic E2)}

WorkflowArena's substrate is unavailable and the paper itself states each
configuration is evaluated once per task without variance
estimates~\cite{workflowarena2026}, so the original criterion (is the author
effect reproducible, and does it vary with difficulty?) is not decidable on
public data. We convert this into an a-priori power analysis. The main-effect MDE
at 145 tasks is dz = 0.234 (below this, undetectable); the difficulty interaction
(72/73 split) has MDE d = 0.468, exactly 2x the main effect---``varies with
difficulty'' is intrinsically twice as hard to detect as ``there is a
difference''; evaluating each task R = 16 times lowers the MDE to 0.059. Monte
Carlo (5000 draws) confirms the H0 rejection rate 0.045 [0.0396, 0.0511] and
power 0.7996 at the MDE. The result is ``cannot-judge is the result'': not that
the author effect is absent, but that the current protocol cannot decide dz $<$
0.234. We learned here that a criterion's direction must be able to fail (the
C3 check originally had the comparison backwards).

\subsection{1.9x internal consistency (AI for science E2)}

The paired human-subject substrate behind the OpenAI 1.9x self-report is
unavailable~\cite{openai-science}. We do an arithmetic self-check on the public
numbers. The five public sub-items have mean 2.142, median 2.05, and geometric
mean 2.112, none equal to 1.9; if 1.9 is an eight-item arithmetic mean, the three
unpublished items would have to average 1.497 (below the lowest public sub-item
1.76); and the feasible mean interval [1.999, 2.400] does not contain 1.9 at all.
The 1.9x headline is therefore not reproducible from public sub-items on the
same caliber, and without pairs or variance the direction and magnitude of any
deviation are mathematically not identifiable---we can only give the time-saving
identity transform (1.9x $\rightarrow$ 47.4\%). We do not claim 1.9x is wrong; we
claim it is not auditable. The caliber lesson: 1.9x (self-reported time ratio)
must never be sorted beside 9.3x (an independently recomputable cost-performance
ratio) from the MIT CRESt result~\cite{crest}.

\subsection{Explanation--energy geometry (foundation-model efficiency E2)}

AROMA reports that the top-4 sensitivity subspace explains 89.82\% of the
second moment yet carries only 9.52\% of median activation
energy~\cite{aroma2026}, which reads like a method finding high-leverage
directions. An isotropic null model (d = 128, rank r = 4, g and h isotropic)
shows this split is the low-rank geometry universal when $r \ll d$: theoretical
median energy r/d = 3.125\%, model null median 2.63\%, and corr(A, E) = 0.011
(sensitivity concentration and energy share are not coupled in the null at all);
the A-null 99th percentile 0.220 matches the paper's own isotropic-null report of
0.219. So 9.52\% is a geometric necessity, neither a defect nor an achievement of
the method; the real increment is that 89.82\% far exceeds its own isotropic
null (10{,}000 rank-matched draws), while the ``high explanation + low energy''
pair carries no methodological information. We note explicitly that we measure
only the A--E coupling geometry, not the +8.00pp count gain, which we did not
re-measure. The null energy uses the analytic $\mathrm{Beta}(r/2,(d-r)/2)$ law
(equivalent and orders of magnitude faster than per-sample eigendecomposition),
and the A-null 99th percentile is sensitive to the number of gradient samples N
(16 $\rightarrow$ 0.3985, 40 $\rightarrow$ 0.2199, 64 $\rightarrow$ 0.1676), so
AROMA's 0.219 corresponds to N $\approx$ 40.

\subsection{92.1\% stratified re-computation (embodied E2)}

The embodied card looked blocked on ``real-robot access''. But its pre-registered
criterion---is the 92.1\% aggregate a few-scenario high mean, and is the tail
stable?---is decidable from ACE-CaP-0's published Table~1 alone~\cite{acecap0}.
The design is balanced: eight scenarios at 30 trials each, so weighted equals
unweighted at 0.9208. The leave-one-out range is 2.38pp; the lowest scenario
(Bread Packing, 83.3\%) has a Wilson CI [66.4\%, 92.7\%] that covers the overall
rate; the between-scenario variance over pure binomial noise ratio is 1.04; and
$\chi^2(7) = 7.259$, p = 0.402, meaning scenario differences are consistent with
binomial sampling noise. The 92.1\% is a flat result, not a mean inflated by a few
strong scenarios. The honest boundary is that this uses only Table~1 aggregate
numbers with no per-trial series, so between-scenario variance is computable but
within-scenario variance is not; with n = 30, ``tail not significantly deviant''
is a low-power non-result, not proof of stability; and it is an arithmetic
consistency check, not a physical reproduction. The GR00T2 closed-loop
experiment (a different, more ambitious test) still awaits artifact release but
does not block this card's own criterion.
\section{Reusable Lessons}
\label{sec:lessons}

The eight diagnostics converge on a small set of portable lessons. We state them
as design rules, each tied to the apparatus failure that motivated it.

\begin{enumerate}
  \item \textbf{Every criterion must carry its own counter-evidence.} A criterion
  that cannot turn red is decoration. Across the eight, the failure modes that
  produced fake-green verdicts were a quantization floor (NSP slope), a
  mis-aligned evaluation arm (multi-view), a template shortcut (DoM probe), a
  cancellation artifact from a correlation that was true for all late steps
  (termination lag), a ``nan'' returned by a non-central-$t$ CDF for large
  non-centrality (workflow power), a reversed criterion direction (1.9x), and a
  slow run that disagreed with a published table because a caliber parameter was
  wrong (AROMA). Each was caught only because we injected a known violation and
  required a red verdict.

  \item \textbf{Fix the criterion before the run, then let it fail.} The NSP
  three-arm is the clearest case: the pre-registered ``lowest slope'' criterion
  failed, and we reported that rather than swapping in the secondary finding as
  the headline. The secondary finding (lowest absolute error) is stated as a
  secondary finding, with the reason the slope criterion is unreliable on stable
  dynamics.

  \item \textbf{Diagnose apparatus defects before changing them.} Each of the
  fixes above was preceded by a concrete diagnosis (the quantization floor, the
  mis-aligned arm, the nan branch), not by a tweak that made the number look
  better. Where we changed a prior from 0.75 to 0.95 in the dynamics, it was a
  physical argument---a strongly stable system attenuates error so much that the
  arms become indistinguishable---made before observing the result, not after.

  \item \textbf{``Substrate unavailable'' is not the same as ``no conclusion''.}
  Power analysis converts ``cannot measure'' into ``MDE = 0.234, R = 16 would
  suffice''; arithmetic self-check converts ``not auditable'' into ``the feasible
  interval excludes the headline''; null geometry converts ``method paradox''
  into ``r $\ll$ d arithmetic''; stratified re-computation converts a
  ``real-robot experiment'' mean into a distributional question answerable from
  a published table. The four unavailable-substrate cards produced four citable
  quantitative statements, not four shrugs.

  \item \textbf{Conclusion scope is fixed by the apparatus, not the model.}
  This is the paper's central claim. The same nominal quantity (0.66) means
  three different things across the embodied literature (a sim-to-real
  reproduction rate, a parts Jaccard index, and an AUC with p = 0.014); 3.1\% and
  $-$7.9\% belong to two different models and cannot be merged into an interval;
  1.9x and 9.3x are different calibers and cannot be ranked. Every number must
  travel with its ruler.

  \item \textbf{Negative and non-result verdicts are first-class.} Four of the
  eight diagnostics return a failed criterion, a mixed result, a cannot-judge,
  or a not-identifiable verdict. We count these as results because they bound
  what can be claimed, and because silently converting them into positive
  findings is exactly the narrative-to-evidence error this paper targets.
\end{enumerate}

Together these rules operationalize the three engineering principles that emerge
from the six directions: \emph{verifiable state prediction} (supervision anchored
to physical or experimental quantities), \emph{auditable failure modes}
(evaluation decomposed to specific failures with variance or intervals), and
\emph{offline generation plus verification} (execution decoupled from checking so
failures can be replayed). Each principle is an instance of the same discipline
the eight diagnostics enact: make the claim checkable, make the check falsifiable,
and state what the check does not cover.
\section{Limits and Citation Discipline}
\label{sec:limits}

\paragraph{Code/data/weight availability.} The full texts of CF-WAM, AROMA,
Termination Lag, and Agentic-OPD were retrieved and checked; none releases code
or data at the time of writing (CF-WAM promises release ``upon acceptance'').
Numbers attributed to those methods are therefore author-reported and cannot be
independently reproduced. The GR00T2 / world-action-model weights were likewise
not released: the NVIDIA/Isaac-GR00T repository's latest tag at check time was
n1.7 (2026-04-18), with no GR00T2/N2/WAM artifact. We deliberately do not build
unfaithful proxies that would masquerade as reproductions; where we used a
synthetic or null-model apparatus, we say so on the same line as the result.

\paragraph{Compute and substrate.} Some diagnostics used a borrowed AMD Radeon
Cloud instance (gfx1100, 48GiB shared) for the synthetic three-arm, multi-view,
and DoM-probe experiments; the rest are pure numpy/standard-library computations
that need no accelerator. Each diagnostic reports its own substrate. This paper
is a methodological audit, not a claim of large-scale reproduction, and the
apparatus scale of each diagnostic bounds the generality of its verdict.

\paragraph{Evidence grading is not truth judgment.} \texttt{ANNOUNCEMENT\_ONLY}
means the current public material does not permit independent recomputation; it
does not assert the result is wrong. Likewise ``not identifiable'' and
``cannot judge'' are statements about the information content of the available
data, not about the underlying phenomenon.

\paragraph{Caliber discipline.} Every arXiv identifier cited here was checked for
number--title agreement against the abstract page and is recorded as VERIFIED in
the project ledger (37 identifiers). All reported percentages are re-derivable
from the corresponding \texttt{evidence/} artifacts. Vendor-reported quantities
(protocol download counts, the 1.9x figure, GeneScientist and NanoTitan numbers)
are explicitly marked as self-reported. Three recalibrations matter for reading
the numbers: RoboBuilder's 66\% is a sim-pass-to-real-pass \emph{reproduction
rate}, not a correlation coefficient; the 3.1\% and $-$7.9\% figures in the
reward-hacking monitoring paper belong to two different models; and any 0.66
citation must name its ruler (reproduction rate vs Jaccard vs AUC vs
correlation).

\paragraph{Reads we do not recommend.} Do not sort the 9.3x independently
recomputable cost-performance ratio beside the 1.9x self-reported time ratio; do
not tabulate the 66\% sim-to-real reproduction rate with the 92.1\% true-robot
success rate; and do not directly compare the long-context PPL numbers of the
three training-free extension routes, which use different bases and different
evaluation protocols.

\paragraph{Reproducibility.} The eight diagnostics, their self-tests, their
landed JSON/text evidence, and the evidence-verification ledger live in the
project tree alongside this manuscript. Every verdict in Table~\ref{tab:eight}
carries the command that reproduces it; the counter-evidence self-tests are part
of the delivered scripts, so a reader can re-run the red cases as well as the
green ones.

\begin{thebibliography}{9}
\bibitem[Ha and Schmidhuber(2018)]{worldmodels2018} Ha, D. and Schmidhuber, J. (2018). World Models. \emph{NeurIPS}. arXiv:1803.10122.
\bibitem[(2025)]{statelaws2025} (2025). Scaling Laws for State Dynamics in Large Language Models. arXiv:2505.14892.
\bibitem[(2026a)]{wam2026} (2026). From World Models to World Action Models: Rethinking Next-State Prediction. arXiv:2609.34414.
\bibitem[(2026b)]{orca2026} (2026). Orca: The World is in Your Mind. arXiv:2606.30534.
\bibitem[(2025b)]{emu35} (2025). Emu3.5: Native Multimodal Models are World Learners. arXiv:2510.26583.
\bibitem[(2026c)]{dom2026} (2026). Monitoring and Discovering Reward Hacking with Internal Representations. arXiv:2609.19101.
\bibitem[(2026d)]{termlag2026} (2026). Termination Lag: Coding Agents Solve Early but Stop Late. arXiv:2611.14960.
\bibitem[(2026e)]{workflowarena2026} (2026). WorkflowArena: Benchmarking Executable Workflow Orchestration in LLM-Based Agents. arXiv:2611.12856.
\bibitem[OpenAI(2026)]{openai-science} OpenAI (2026). Science at the Frontier: How scientists use ChatGPT. Report, 2026-11-16.
\bibitem[(2025c)]{crest} (2025). Nature 647:390--396. doi:10.1038/s41586-025-09640-5.
\bibitem[(2026f)]{aroma2026} (2026). AROMA: Adaptive Representation-Oriented Multimodal Actuation for Vision-Language Models. arXiv:2611.15041.
\bibitem[(2026g)]{acecap0} (2026). ACE-CaP-0: Real-World Code-as-Policies in the Era of Experience. arXiv:2611.16872.
\bibitem[(2026h)]{agenticopd} (2026). Agentic-OPD: Recovering Forgotten Agent Capabilities via On-Policy Distillation. arXiv:2611.12799.
\bibitem[(2026i)]{misdiag} (2026). Aggregate Accuracy Hides Diagnosis-Specific Failure in Medical Vision-Language Models. arXiv:2611.14975.
\end{thebibliography}

\end{document}